% Part: first-order-logic
% Chapter: sequent-calculus
% Section: proving-things-quant

\documentclass[../../../include/open-logic-section]{subfiles}

\begin{document}

\olfileid[tr]{fol}{seq}{prq}

\olsection{Niceleyiciler İçeren \usetoken{P}{derivation}}

\begin{ex}
$\lexists[x][\lnot !A(x)] \Sequent \lnot \lforall[x][!A(x)]$ sekantının
bir $\Log{LK}$-!!{derivation}{}ini verin.

Niceleyicilerle çalışırken özdeğişken koşulunu ihlal etmemeye dikkat etmemiz
gerekir ve bu bazen belirli çıkarımları uygulama sıramızla oynamamızı
gerektirir. Genel olarak özdeğişken koşuluna bağlı kuralları önce ele almaya
çalışmak yararlıdır (tamamlanmış ispatta bunlar daha aşağıda bulunur). Ayrıca
ileriye bakıp başlangıç sekantının nasıl görüneceğini kestirmeye çalışmak da
iyi bir fikirdir. Bizim örneğimizde başlangıç sekantı $!A(a) \Sequent !A(a)$
gibi bir şey olmalıdır. Bu, niceleyici kurallarını ``tersine'' uygularken hem
$\lforall$ hem de $\lexists$ kuralı için aynı terimi---$a$ diyeceğimiz
terimi---seçmemiz gerektiği anlamına gelir. Her kural için farklı bir terim
seçseydik $!A(a) \Sequent !A(b)$ gibi, elbette türetilemeyen bir sekant elde
ederdik.

Her zamanki gibi şu ifadeyi yazarak başlayalım:
\begin{prooftree}
\AxiomC{}
\UnaryInf$\lexists[x][\lnot !A(x)] \fCenter \lnot \lforall[x][!A(x)]$
\end{prooftree}
\LeftR{\exists} veya \RightR{\lnot} kuralını uygulayabiliriz.
\LeftR{\exists} kuralı özdeğişken koşuluna bağlı olduğundan onu geciktirmeden
ele almak iyi bir fikirdir; dolayısıyla önce bu kuralı
uygulayalım.
\begin{prooftree}
\AxiomC{}
\UnaryInf$ \lnot !A(a) \fCenter \lnot \lforall[x][!A(x)]$
\RightLabel{\LeftR{\lexists}}
\UnaryInf$ \lexists[x][\lnot !A(x)] \fCenter \lnot \lforall[x][!A(x)]$
\end{prooftree}
\LeftR{\lnot} ve \RightR{\lnot} kurallarını geriye doğru uygulayınca
şunu elde ederiz:
\begin{prooftree}
\AxiomC{}
\UnaryInf$\lforall[x][!A(x)] \fCenter !A(a)$
\RightLabel{\LeftR{\lnot}}
\UnaryInf$\lnot !A(a), \lforall[x][!A(x)] \fCenter $
\RightLabel{\LeftR{\Exchange}}
\UnaryInf$\lforall[x][!A(x)], \lnot !A(a) \fCenter $
\RightLabel{\RightR{\lnot}}
\UnaryInf$ \lnot !A(a) \fCenter \lnot \lforall[x] !A(x)$
\RightLabel{\LeftR{\lexists}}
\UnaryInf$ \lexists[x] \lnot !A(x) \fCenter \lnot \lforall[x] !A(x)$
\end{prooftree}
Bu noktada tek seçeneğimiz \LeftR{\forall} kuralını uygulamaktır. Bu kural
özdeğişken kısıtlamasına bağlı olmadığından sorun yoktur. Bir başlangıç sekantı
($!A(a) \Sequent !A(a)$ biçiminde) elde etmeye çalıştığımızı hatırlayın;
dolayısıyla kuralı uygularken $!A$'nın argümanı olarak $a$'yı seçmeliyiz.
\begin{prooftree}
\Axiom$!A(a) \fCenter !A(a)$
\RightLabel{\LeftR{\lforall}}
\UnaryInf$\lforall[x][!A(x)] \fCenter !A(a)$
\RightLabel{\LeftR{\lnot}}
\UnaryInf$\lnot !A(a), \lforall[x][!A(x)] \fCenter $
\RightLabel{\LeftR{\Exchange}}
\UnaryInf$\lforall[x][!A(x)], \lnot !A(a) \fCenter $
\RightLabel{\RightR{\lnot}}
\UnaryInf$ \lnot !A(a) \fCenter \lnot \lforall[x][!A(x)]$
\RightLabel{\LeftR{\lexists}}
\UnaryInf$ \lexists[x][ \lnot !A(x)] \fCenter \lnot \lforall[x][!A(x)]$
\end{prooftree}
Özellikle niceleyicilerle çalışırken özdeğişken koşulunun ihlal edilmediğini
bu noktada iki kez denetlemek önemlidir. Özdeğişken koşuluna bağlı olarak
uyguladığımız tek kural \LeftR{\exists} olduğundan ve özdeğişken $a$ bu
kuralın alt sekantında (son sekantta) bulunmadığından bu doğru bir
!!{derivation}{}dir.
\end{ex}

\begin{prob}
Aşağıdaki sekantlar için !!{derivation}s verin:
\begin{enumerate}
\item $\Sequent (\lforall[x][!A(x)] \land \lforall[y][!B(y)]) \lif
\lforall[z][(!A(z) \land !B(z))]$.
\item $\Sequent (\lexists[x][!A(x)] \lor \lexists[y][!B(y)]) \lif
\lexists[z][(!A(z) \lor !B(z))]$.
\item $\lforall[x][(!A(x) \lif !B)] \Sequent \lexists[y][!A(y)] \lif !B$.
\item $\lforall[x][\lnot !A(x)] \Sequent \lnot\lexists[x][!A(x)]$.
\item $\Sequent \lnot\lexists[x][!A(x)] \lif \lforall[x][\lnot !A(x)]$.
\item $\Sequent \lnot\lexists[x][\lforall[y][((!A(x,y) \lif \lnot
!A(y,y)) \land (\lnot !A(y,y) \lif !A(x,y)))]]$.
\end{enumerate}
\end{prob}

\begin{prob}
Aşağıdaki sekantlar için !!{derivation}s verin:
\begin{enumerate}
\item $\Sequent \lnot\lforall[x][!A(x)] \lif \lexists[x][\lnot!A(x)]$.
\item $(\lforall[x][!A(x)] \lif !B) \Sequent \lexists[y][(!A(y) \lif !B)]$.
\item $\Sequent \lexists[x][(!A(x) \lif \lforall[y][!A(y)])]$.
\end{enumerate}
(Bunların tümü $\RightR{\Contraction}$ kuralını gerektirir.)
\end{prob}

\end{document}
